\subsection{Case study: a one-loop Majorana-neutrino and dark-matter model}
\label{case:2608-12646}

\paragraph{Paper identity and theme.}
The source is Mohamed Belfkir, Mohamed Amin Loualidi, and Salah Nasri,
``A Novel One-loop Model for Majorana Neutrino Mass and Dark Matter,''
arXiv:2608.12646v1. The paper constructs a $T4$-$3$-$i$ realization in
which a Dirac mediator removes a lower-order seesaw contribution, a separate
Majorana fermion supplies lepton-number violation, and an exact $Z_2$ symmetry
stabilizes dark matter. It derives scalar, neutrino, charged-lepton-flavor,
and dark-matter observables and fits four combinations of mass ordering and
dark-matter identity. Claims attributed below to ``the paper'' are source
statements; comments headed ``Interface analysis'' are the present stress-test
assessment. The machine-readable inventory contains 53 consolidated candidates, including one deliberate composite. This is candidate-stage stress data: promotion to an accepted ScientificClaim still requires the later atomicity gate.

\paragraph{Source census method.}
The census read all of
\path{tmp/claim-interface-study/extracted/2608.12646/main.tex}, including
macros, tables, figure captions, and the complete unitarity appendix. A
claim-bearing unit is a prose paragraph, list item, equation-plus-lead, table,
or caption that contributes a definition, assumption, imported result,
derivation, analytical result, numerical finding, limitation, or
interpretation. Headings, layout, acknowledgements, and bibliography commands
were excluded. Abstract, introduction, body, and conclusion repetitions were
counted as occurrences and consolidated only after their scopes were compared.

\begin{center}
\begin{tabular}{lrr}
\hline
Source division & Claim-bearing units & Consolidated candidates \\
\hline
Abstract and introduction & 10 & 6 \\
Topology and charge classification & 19 & 10 \\
Scalar potential and spectrum & 18 & 6 \\
Neutrino mass and cLFV & 24 & 7 \\
Theoretical and experimental constraints & 13 & 5 \\
Dark-matter constraints & 17 & 8 \\
Statistical framework & 9 & 2 \\
Fit results and benchmarks & 29 & 13 \\
Summary, limitations, and outlook & 5 & 3 \\
Unitarity appendix & 7 & 1 \\
\hline
Total occurrences & 151 & 53 unique candidates \\
\hline
\end{tabular}
\end{center}

Some candidates use support from more than one division, and some dense source
paragraphs split into several independent predications; consequently the
candidate total is not an additive partition of the right column.

\paragraph{Compact claim-family census.}
\begin{center}
\begin{tabular}{p{0.23\linewidth}p{0.50\linewidth}r}
\hline
Family & Examples retained & Candidates \\
\hline
Background, novelty, and model scope & Genuineness criterion, priority claim,
minimal field content & 5 \\
Classification and setup & Four charge families, allowed $\alpha$ variants,
$Z_2$ and lepton-number design & 8 \\
Analytical field-theory results & Scalar matrices, one-loop neutrino mass,
cLFV, unitarity, Higgs and oblique observables & 18 \\
Dark-matter assumptions and mechanisms & Freeze-out, coannihilation, tree- and
loop-level direct detection & 8 \\
Statistical and numerical findings & Likelihood, acceptance counts, fitted
ranges, benchmark observables, viability & 11 \\
Limitations and interpretation & Collider omission, future work, region labels
& 3 \\
\hline
\end{tabular}
\end{center}

\paragraph{Representative difficult splits.}
The scalar-spectrum candidate retains one deliberate bundle for identifier stability, but its conditions are clause-specific. The chosen inert-vacuum statement says that only the neutral CP-even Standard Model Higgs component obtains $v=246\,\mathrm{GeV}$ while $\phi$ and $\phi^{\prime}$ remain inert. Separately, taking the parameters of $\mathcal V_B$ real avoids explicit CP violation, and the $\mu_\phi$ term then mixes the neutral fields within distinct CP-even and CP-odd sectors. The real-parameter condition must not be attached to the VEV proposition. Accepted promotion should split the VEV choice, the no-explicit-CP assumption, and the mixing/spectrum consequence.

The model-building argument is not one indivisible claim. The inventory
separates the conditional obstruction---a Majorana singlet or triplet mediator
permits a lower-order seesaw---from the authors' construction using Dirac $N$,
Majorana $\psi$, and exact $Z_2$. It then separates the resulting field
inventory from the later numerical viability claim. This prevents the
existence of a neutral odd state from being conflated with successful relic
abundance and direct detection.

The most compact analytical result is
\[
\begin{aligned}
 M_\nu &= \Lambda\left(Yy^T+yY^T\right),\\
 \Lambda &=-\mu_\phi\langle H^0\rangle^2
 \frac{M_\psi}{M_N}y'_{11}
 I_3(m_\phi^2,m_{\phi'}^2,M_\psi^2).
\end{aligned}
\]
For one $N$ and one $\psi$,
\[
\begin{aligned}
 \operatorname{rank}(M_\nu)&\leq 2,\\
 \det M_\nu&=0
 \quad\text{for generic rank two},
\end{aligned}
\]
so one neutrino is massless at this order. The record keeps ``generic flavor
vectors'' and ``at this order'' as scope. The normal- and inverted-ordering
mass intervals from the scan are distinct numerical claims rather than part of
the rank theorem.

A similar split is needed for direct detection. For fermionic dark matter the
paper derives, in the heavy-charged-scalar limit,
\[
 g_{h\psi\psi}^{\mathrm{eff}}
 \simeq-\frac{\kappa_1v}{32\pi^2}
 \left(\sum_\alpha|y_\alpha|^2\right)
 \frac{M_\psi}{m_{\phi^\pm}^2}.
\]
This decoupling formula is an analytical candidate. The statement that all
accepted fermionic points lie below xenon limits is separate because it also
depends on scan ranges, cuts, and external exclusion curves. For scalar dark
matter, the tree-level $hH_1^0H_1^0$ coupling is likewise distinct from the
reported accepted cross-section range. The accepted $H_1^0$ mass range starts
near $62.8\,\mathrm{GeV}$ and reaches about $896\,\mathrm{GeV}$ for normal
ordering and $773\,\mathrm{GeV}$ for inverted ordering; the medians are near
$70\,\mathrm{GeV}$, and the respective best fits are dominated by
$H_1^0A_1^0$ and $H_1^0A_2^0$ coannihilation.

The relic-density discussion combines a standard equation with an
interpretive region label:
\[
\begin{aligned}
 \frac{dY}{dx}
 &=-\frac{s}{Hx}\langle\sigma_{\mathrm{eff}}v\rangle
 \left(Y^2-Y_{\mathrm{eq}}^2\right),\\
 r_i(x)&=\frac{g_i(1+\Delta_i)^{3/2}e^{-x\Delta_i}}
 {\sum_k g_k(1+\Delta_k)^{3/2}e^{-x\Delta_k}}.
\end{aligned}
\]
The equations and radiation-dominated freeze-out assumption normalize. The
claim that a viable sample is best called a ``Higgs-funnel/coannihilation
region'' does not normalize without an explicit region-membership criterion.

\paragraph{Source-fidelity decisions.}
All JSONL spans are one-indexed complete-line ranges whose text is copied
byte-for-byte. Tables were retained when they define finite mappings, scan
ranges, acceptance counts, or benchmarks. Figure captions alone were not
used to infer trends absent from the prose. The priority statement remains
hedged as ``to the authors' knowledge'' and is attributed to the paper rather
than accepted as a verified literature fact. Imported likelihood profiles,
experimental contours, and projected sensitivities remain external evidence.
Approximations such as $\Delta U=0$, the Standard-Model-like Higgs total width,
leading-chiral conversion formulas, and a radiation-dominated thermal history
are attached to the claims that depend on them.

\paragraph{Interface analysis: MathClaimIR boundaries.}
The 53 candidates contain explicit equations and finite mappings that normalize
well; model-wide causal statements and scan-conditioned findings generally
normalize only partially. Priority, minimality, qualitative region labels,
and research-scope statements do not have faithful MathClaimIR forms. In
particular, the existential claim that all four scenarios are viable is not an
unqualified theorem: it means that each finite scan contains at least one point
passing the stated implementation, priors, ranges, likelihoods, and cuts.

\paragraph{Interface pressures.}
The case exposes nine pressures: conditional ``genuine topology'' status;
priority claims requiring external verification; physical exclusions versus
authors' convenience restrictions in a classification; assumptions propagated
across long derivations; removable singular domains in reconstructed
parameters; mixed likelihood and hard-cut semantics; selection-induced
conclusions; benchmark evidence versus population-level claims; and large,
ordered matrix collections in the appendix. It also shows why an external
experimental curve cannot be reduced to the prose phrase ``below the bound''.
These pressures concern evidence, scope, and provenance and are recorded
without changing the proposed interface schema.

\paragraph{Provisional verdict.}
The interface is strong for explicit Lagrangians, matrices, inequalities,
finite classifications, and parameter-to-observable mappings. It remains
usable for phenomenology only if assumptions and scan provenance stay attached
to each candidate. A normalization refusal is essential for priority,
minimality, qualitative region, and future-work claims. The case therefore supports the five-field accepted identity kernel with exact
source spans and an atomic, source-faithful statement. Scope, modality, and
attribution remain mandatory candidate and admission semantics that must be
projected into that statement, while optional mathematical normalization and
visible pressure flags stay outside accepted identity. No physics-specific
expansion of the Registry is required.